\chapter{2021年全国甲卷（文）}

\section{单选题}

\begin{problem}
设集合 \(M=\{1,3,5,7,9\}\)，\(N=\{x\mid 2x>7\}\)，则 \(M\cap N=\)（\quad）

\choices
    {\(\{7,9\}\)}
    {\(\{5,7,9\}\)}
    {\(\{3,5,7,9\}\)}
    {\(\{1,3,5,7,9\}\)}

\answer{B}
\end{problem}

\begin{problem}
为了解某地农村经济情况，对该地农户家庭年收入进行抽样调查，将农户家庭年收入的调查数据整理得到如下频率分布直方图：

\bitmapfigure[max width=0.90\linewidth]{img/2021/national_paper_a_liberal/q2_fig1.png}

根据此频率分布直方图，下面结论中不正确的是（\quad）

\choices
    {该地农户家庭年收入低于 \(4.5\) 万元的农户比率估计为 \(6\%\)}
    {该地农户家庭年收入不低于 \(10.5\) 万元的农户比率估计为 \(10\%\)}
    {估计该地农户家庭年收入的平均值不超过 \(6.5\) 万元}
    {估计该地有一半以上的农户，其家庭年收入介于 \(4.5\) 万元至 \(8.5\) 万元之间}

\answer{C}
\end{problem}

\begin{problem}
已知 \((1-\i)^2z=3+2\i\)，则 \(z=\)（\quad）

\choices
    {\(-1-\frac{3}{2}\i\)}
    {\(-1+\frac{3}{2}\i\)}
    {\(-\frac{3}{2}+\i\)}
    {\(-\frac{3}{2}-\i\)}

\answer{B}
\end{problem}

\begin{problem}
下列函数中是增函数的是（\quad）

\choices
    {\(f(x)=-x\)}
    {\(f(x)=\left(\frac{2}{3}\right)^x\)}
    {\(f(x)=x^2\)}
    {\(f(x)=\sqrt[3]{x}\)}

\answer{D}
\end{problem}

\begin{problem}
点 \((3,0)\) 到双曲线 \(\frac{x^2}{16}-\frac{y^2}{9}=1\) 的一条渐近线的距离为（\quad）

\choices
    {\(\frac{9}{5}\)}
    {\(\frac{8}{5}\)}
    {\(\frac{6}{5}\)}
    {\(\frac{4}{5}\)}

\answer{A}
\end{problem}

\begin{problem}
青少年视力是社会普遍关注的问题，视力情况可借助视力表测量. 通常用五分记录法和小数记录法记录视力数据，五分记录法的数据 \(L\) 和小数记录法的数据 \(V\) 满足 \(L=5+\lg V\). 已知某同学视力的五分记录法的数据为 \(4.9\)，则其视力的小数记录法的数据约为（\quad）（\(\sqrt[10]{10}\approx 1.259\)）

\choices
    {\(1.5\)}
    {\(1.2\)}
    {\(0.8\)}
    {\(0.6\)}

\answer{C}
\end{problem}

\begin{problem}
在一个正方体中，过顶点 \(A\) 的三条棱的中点分别为 \(E\)，\(F\)，\(G\). 该正方体截去三棱锥 \(A-EFG\) 后，所得多面体的三视图中，正视图如图所示，则相应的侧视图是（\quad）

\FigureLayoutDeclare{img/2021/national_paper_a_liberal/q7_fig1.png}{4.0cm}{31bp 9bp 30bp 18bp}
\bitmapfigure[max width=7cm]{img/2021/national_paper_a_liberal/q7_fig1.png}

\choices
    {\choicebitmap[width=0.15\paperwidth]{img/2021/national_paper_a_liberal/q7_optA.png}}
    {\choicebitmap[width=0.15\paperwidth]{img/2021/national_paper_a_liberal/q7_optB.png}}
    {\choicebitmap[width=0.15\paperwidth]{img/2021/national_paper_a_liberal/q7_optC.png}}
    {\choicebitmap[width=0.15\paperwidth]{img/2021/national_paper_a_liberal/q7_optD.png}}

\answer{D}
\end{problem}

\begin{problem}
在 \(\triangle ABC\) 中，已知 \(B=120^\circ\)，\(AC=\sqrt{19}\)，\(AB=2\)，则 \(BC=\)（\quad）

\choices
    {\(1\)}
    {\(\sqrt{2}\)}
    {\(\sqrt{5}\)}
    {\(3\)}

\answer{D}
\end{problem}

\begin{problem}
记 \(S_n\) 为等比数列 \(\{a_n\}\) 的前 \(n\) 项和. 若 \(S_2=4\)，\(S_4=6\)，则 \(S_6=\)（\quad）

\choices
    {7}
    {8}
    {9}
    {10}

\answer{A}
\end{problem}

\begin{problem}
将 \(3\) 个 \(1\) 和 \(2\) 个 \(0\) 随机排成一行，则 \(2\) 个 \(0\) 不相邻的概率为（\quad）

\choices
    {0.3}
    {0.5}
    {0.6}
    {0.8}

\answer{C}
\end{problem}

\begin{problem}
若 \(\alpha\in\left(0,\frac{\pi}{2}\right)\)，\(\tan 2\alpha=\frac{\cos\alpha}{2-\sin\alpha}\)，则 \(\tan\alpha=\)（\quad）

\choices
    {\(\frac{\sqrt{15}}{15}\)}
    {\(\frac{\sqrt{5}}{5}\)}
    {\(\frac{\sqrt{5}}{3}\)}
    {\(\frac{\sqrt{15}}{3}\)}

\answer{A}
\end{problem}

\begin{problem}
设 \(f(x)\) 是定义域为 \(\R\) 的奇函数，且 \(f(1+x)=f(-x)\). 若 \(f\left(-\frac{1}{3}\right)=\frac{1}{3}\)，则 \(f\left(\frac{5}{3}\right)=\)（\quad）

\choices
    {\(-\frac{5}{3}\)}
    {\(-\frac{1}{3}\)}
    {\(\frac{1}{3}\)}
    {\(\frac{5}{3}\)}

\answer{C}
\end{problem}

\section{填空题}

\begin{problem}
若向量 \(\bs a\)，\(\bs b\) 满足 \(|\bs a|=3\)，\(|\bs a-\bs b|=5\)，\(\bs a\cdot\bs b=1\)，则 \(|\bs b|=\fillinblank{}\).

\answer{\(3\sqrt{2}\)}
\end{problem}

\begin{problem}
已知一个圆锥的底面半径为 \(6\)，其体积为 \(30\pi\)，则该圆锥的侧面积为\fillinblank{}.

\answer{\(39\pi\)}
\end{problem}

\begin{problem}
已知函数 \(f(x)=2\cos(\omega x+\varphi)\) 的部分图像如图所示，则 \(f\left(\frac{\pi}{2}\right)=\fillinblank{}\).

\FigureLayoutDeclare{img/2021/national_paper_a_liberal/q15_fig1.png}{12.0cm}{0bp 5bp 35bp 19bp}
\bitmapfigure[max width=6.0cm]{img/2021/national_paper_a_liberal/q15_fig1.png}

\answer{\(-\sqrt{3}\)}
\end{problem}

\begin{problem}
已知 \(F_1\)，\(F_2\) 为椭圆 \(C:\frac{x^2}{16}+\frac{y^2}{4}=1\) 的两个焦点，\(P\)，\(Q\) 为 \(C\) 上关于坐标原点对称的两点，且 \(|PQ|=|F_1F_2|\)，则四边形 \(PF_1QF_2\) 的面积为\fillinblank{}.

\answer{8}
\end{problem}

\section{解答题}

\begin{problem}
甲，乙两台机床生产同种产品，产品按质量分为一级品和二级品. 为了比较两台机床产品的质量，分别用两台机床各生产了 \(200\) 件产品，产品的质量情况统计如下表：

\begin{center}
\begin{tabular}{|c|c|c|c|}
\hline
 & 一级品 & 二级品 & 合计 \\
\hline
甲机床 & 150 & 50 & 200 \\
\hline
乙机床 & 120 & 80 & 200 \\
\hline
合计 & 270 & 130 & 400 \\
\hline
\end{tabular}
\end{center}

\begin{enumerate}
\item 甲机床，乙机床生产的产品中一级品的频率分别是多少？
\item 能否有 \(99\%\) 的把握认为甲机床的产品质量与乙机床的产品质量有差异？
\end{enumerate}

附：\(K^2=\dfrac{n(ad-bc)^2}{(a+b)(c+d)(a+c)(b+d)}\)

\begin{center}
\begin{tabular}{|c|c|c|c|}
\hline
\(P(K^2\ge k)\) & 0.050 & 0.010 & 0.001 \\
\hline
\(k\) & 3.841 & 6.635 & 10.828 \\
\hline
\end{tabular}
\end{center}
\end{problem}

\solution{短答案：甲机床、乙机床生产的一级品频率分别为 \(\frac{3}{4}\)、\(\frac{3}{5}\)；\(K^2\approx10.256>6.635\)，有 \(99\%\) 的把握认为两台机床的产品质量有差异.}

\begin{problem}
记 \(S_n\) 为数列 \(\{a_n\}\) 的前 \(n\) 项和，已知 \(a_n>0\)，\(a_2=3a_1\)，且数列 \(\{\sqrt{S_n}\}\) 是等差数列，证明：\(\{a_n\}\) 是等差数列.
\end{problem}

\solution{短答案：\(a_n=(2n-1)a_1\)，故 \(\{a_n\}\) 是等差数列.}

\begin{problem}
已知直三棱柱 \(ABC-A_1B_1C_1\) 中，侧面 \(AA_1B_1B\) 为正方形，\(AB=BC=2\)，\(E\)，\(F\) 分别为 \(AC\) 和 \(CC_1\) 的中点，\(BF\perp A_1B_1\).

\begin{enumerate}
\item 求三棱锥 \(F-EBC\) 的体积；
\item 已知 \(D\) 为棱 \(A_1B_1\) 上的点，证明：\(BF\perp DE\).
\end{enumerate}

\bitmapfigure[max width=5.0cm]{img/2021/national_paper_a_liberal/q19_fig1.png}
\end{problem}

\solution{短答案：第一问体积为 \(\frac{1}{3}\)；第二问 \(BF\perp DE\).}

\begin{problem}
设函数 \(f(x)=a^2x^2+ax-3\ln x+1\)，其中 \(a>0\).

\begin{enumerate}
\item 讨论 \(f(x)\) 的单调性；
\item 若 \(y=f(x)\) 的图象与 \(x\) 轴没有公共点，求 \(a\) 的取值范围.
\end{enumerate}
\end{problem}

\solution{短答案：\(f(x)\) 在 \(\left(0,\frac{1}{a}\right)\) 上单调递减，在 \(\left(\frac{1}{a},+\infty\right)\) 上单调递增；\(a\in\left(\frac{1}{\e},+\infty\right)\).}

\begin{problem}
抛物线 \(C\) 的顶点为坐标原点 \(O\)，焦点在 \(x\) 轴上，直线 \(l:x=1\) 交 \(C\) 于 \(P\)，\(Q\) 两点，且 \(OP\perp OQ\). 已知点 \(M(2,0)\)，且 \(\odot M\) 与 \(l\) 相切.

\begin{enumerate}
\item 求 \(C\)，\(\odot M\) 的方程；
\item 设 \(A_1\)，\(A_2\)，\(A_3\) 是 \(C\) 上的三个点，直线 \(A_1A_2\)，\(A_1A_3\) 均与 \(\odot M\) 相切. 判断直线 \(A_2A_3\) 与 \(\odot M\) 的位置关系，并说明理由.
\end{enumerate}
\end{problem}

\solution{短答案：\(C:y^2=x\)，\(\odot M:(x-2)^2+y^2=1\)；直线 \(A_2A_3\) 与 \(\odot M\) 相切.}

\section{选考题：解答题}

\begin{problem}
在直角坐标系 \(xOy\) 中，以坐标原点为极点，\(x\) 轴正半轴为极轴建立极坐标系，曲线 \(C\) 的极坐标方程为 \(\rho=2\sqrt{2}\cos\theta\).

\begin{enumerate}
\item 将 \(C\) 的极坐标方程化为直角坐标方程；
\item 设点 \(A\) 的直角坐标为 \((1,0)\)，\(M\) 为 \(C\) 上的动点，点 \(P\) 满足 \(\overrightarrow{AP}=\sqrt{2}\overrightarrow{AM}\)，写出 \(P\) 的轨迹 \(C_1\) 的参数方程，并判断 \(C\) 与 \(C_1\) 是否有公共点.
\end{enumerate}
\end{problem}

\solution{短答案：\(C:(x-\sqrt{2})^2+y^2=2\)；\(C_1\) 的参数方程为
\(\begin{cases}
x=3-\sqrt{2}+2\cos\theta,\\
y=2\sin\theta,
\end{cases}\)
两曲线没有公共点.}

\begin{problem}
已知函数 \(f(x)=|x-2|\)，\(g(x)=|2x+3|-|2x-1|\).

\begin{enumerate}
\item 画出 \(y=f(x)\) 和 \(y=g(x)\) 的图像；
\item 若 \(f(x+a)\ge g(x)\)，求 \(a\) 的取值范围.
\end{enumerate}

\bitmapfigure[float=inline,max width=7.0cm]{img/2021/national_paper_a_liberal/q23_fig1.png}
\FloatBarrier
\end{problem}

\solution{短答案：
\(f(x)=\begin{cases}
2-x,&x<2,\\
x-2,&x\geqslant2,
\end{cases}\quad
g(x)=\begin{cases}
-4,&x< -\frac{3}{2},\\
4x+2,&-\frac{3}{2}\leqslant x<\frac{1}{2},\\
4,&x\geqslant\frac{1}{2},
\end{cases}\)
；\(a\in\left[\frac{11}{2},+\infty\right)\).}
